% ============================================================
%  07_thao_luan.tex  --  Thao luan va ket luan
% ============================================================

\section{Thảo luận và kết luận}
\label{sec:ketluan}

% ────────────────────────────────────────────────────────────
\subsection{Tổng hợp kết quả thực nghiệm}

Bài thực hành đã đạt được đầy đủ các mục tiêu đề ra. Bảng dưới tóm tắt toàn
bộ kết quả:

\begin{table}[H]
\centering
\begin{tabularx}{\textwidth}{|p{4cm}|X|p{2.5cm}|}
\hline
\textbf{Mục tiêu} & \textbf{Kết quả đạt được} & \textbf{Trạng thái} \\ \hline
Phát hiện 3/3 lỗi trên bản vulnerable & ASan báo đủ heap-buffer-overflow, heap-use-after-free, detected memory leaks & Đạt \\ \hline
Xác nhận 0/3 lỗi trên bản secure & Ba kịch bản chạy sạch, exit code 0, không có ERROR & Đạt \\ \hline
8/8 test PASS & Toàn bộ bộ kiểm thử tự động thông qua & Đạt \\ \hline
Đánh giá rủi ro 5 lỗ hổng & Kết quả khớp baseline $\leq 0.05$ điểm & Đạt \\ \hline
Giải thích giới hạn ASan & Thực nghiệm tràn trong struct đã thực hiện & Đạt \\ \hline
\end{tabularx}
\caption{Tổng hợp kết quả so với mục tiêu}
\end{table}

% ────────────────────────────────────────────────────────────
\subsection{Phân tích sâu: vì sao lỗi bộ nhớ vẫn phổ biến?}

Mặc dù các lỗi như CWE-121 đã được biết đến từ thập niên 1980 và có các công
cụ phát hiện tốt, chúng vẫn xuất hiện thường xuyên trong phần mềm thực tế.
Có ba nguyên nhân cốt lõi:

\begin{enumerate}
    \item API kế thừa không an toàn: C không loại bỏ \textit{strcpy},
          \textit{gets}, \textit{sprintf} khỏi thư viện chuẩn. Code cũ dùng
          các hàm này vẫn được biên dịch mà không có cảnh báo bắt buộc (trừ
          khi bật \textit{-Wall -Wextra}).
    \item Kiểm thử không đủ kịch bản: Lỗi chỉ xuất hiện khi đầu vào vượt ngưỡng
          nhất định. Kiểm thử thông thường dùng đầu vào hợp lệ không bao giờ
          kích hoạt lỗi, dẫn đến cảm giác sai về độ an toàn.
    \item Giới hạn công cụ: ASan phát hiện lỗi tại runtime chỉ trên kịch bản
          được chạy. Tràn trong struct, lỗi trong nhánh code ít chạy, hay điều
          kiện race condition có thể không bị bắt.
\end{enumerate}

% ────────────────────────────────────────────────────────────
\subsection{Ý nghĩa của mô hình đánh giá rủi ro}

Kết quả CVSS-lite cho thấy không phải lỗi nào cũng nguy hiểm như nhau. Cùng
là lỗi bộ nhớ, nhưng:

\begin{itemize}
    \item Buffer Overflow (9.8) và Format String (9.8) nguy hiểm gần gấp 4 lần
          NULL Pointer Dereference (2.5) vì chúng có thể dẫn đến thực thi mã
          từ xa không cần đặc quyền.
    \item Use-After-Free (6.9) có Impact tương đương Buffer Overflow
          (Impact = 5.87) nhưng Exploitability thấp hơn nhiều (1.05 so với 3.89)
          do cần tấn công phức tạp hơn.
    \item Memory Leak (5.4) chỉ ảnh hưởng Availability -- không rò rỉ dữ liệu
          hay cho phép thực thi mã -- nên được ưu tiên thấp hơn.
\end{itemize}

Điều này trực tiếp hướng dẫn quyết định phân bổ nguồn lực: vá VULN-001 trước
vì chi phí khai thác thấp nhất và hậu quả nghiêm trọng nhất.

% ────────────────────────────────────────────────────────────
\subsection{Hạn chế của bài thực hành}

\begin{itemize}
    \item Môi trường đơn luồng: Module Inventory không có race condition; lỗi
          Use-After-Free trong môi trường đa luồng còn phức tạp và nguy hiểm
          hơn nhiều (TOCTOU, concurrent free).
    \item Không có fuzzing: Bộ test chỉ kiểm tra ba kịch bản cố định. Fuzzing
          (AFL, libFuzzer) có thể tìm thêm lỗi ở các edge case không được nghĩ
          đến.
    \item Mô hình CVSS-lite đơn giản: Không tính đến Scope (Unchanged/Changed),
          môi trường triển khai cụ thể, hay các biện pháp kiểm soát bổ sung
          (WAF, ASLR, stack canary). Điểm thực tế có thể thấp hơn nếu hệ thống
          có nhiều lớp phòng thủ.
    \item Phụ thuộc vào ASan: Trên Windows, biên dịch và chạy với
          \textit{-fsanitize=address} đòi hỏi môi trường đặc biệt (MinGW,
          Clang-cl). Bài thực hành giả định môi trường Linux/WSL.
\end{itemize}

% ────────────────────────────────────────────────────────────
\subsection{Bài học thực tiễn}

Qua bài thực hành, có thể rút ra các bài học sau:

\begin{enumerate}
    \item ``Secure by default'': Mặc định dùng API an toàn. Thay vì ghi nhớ
          ``đừng dùng \textit{strcpy}'', hãy thiết lập quy tắc
          \textit{-Wformat-security}, \textit{-D\_FORTIFY\_SOURCE=2} trong
          pipeline biên dịch để compiler tự cảnh báo.
    \item Defense in depth: Không có công cụ nào đủ một mình. ASan + code review
          + phân tích tĩnh + fuzzing cộng lại mới tạo ra độ bảo phủ đủ tốt.
    \item Đánh giá rủi ro trước khi vá: Khi có nhiều lỗi, vá tất cả cùng lúc
          không phải lúc nào cũng khả thi. Điểm CVSS giúp quyết định lỗi nào
          cần vá ngay hôm nay, lỗi nào có thể đưa vào sprint tiếp theo.
    \item Kiểm thử sau khi vá: Bản vá đúng không chỉ là ``chạy không báo lỗi
          ASan'' -- bộ kiểm thử tự động xác nhận hành vi đúng và tính nhất quán
          của mô hình rủi ro là lớp xác minh thứ hai quan trọng.
\end{enumerate}

% ────────────────────────────────────────────────────────────
\subsection{Kết luận}

Bài thực hành 03 đã thực hiện đầy đủ chu trình: tái hiện lỗi $\to$
phát hiện bằng công cụ $\to$ hiểu cơ chế $\to$ sửa lỗi $\to$
xác nhận sửa đúng $\to$ đánh giá rủi ro và ưu tiên.

Ba lỗi bộ nhớ kinh điển CWE-121, CWE-416 và CWE-401 đã được tái hiện và phát
hiện bằng AddressSanitizer, sửa theo nguyên tắc lập trình C an toàn, và xác
nhận bằng bộ kiểm thử tự động 8/8 PASS. Mô hình CVSS-lite định lượng hóa rủi
ro, xếp Buffer Overflow là ưu tiên khắc phục cao nhất (9.8/CRITICAL) -- phù
hợp với thực tiễn bảo mật.

Điểm sư phạm quan trọng nhất: ASan không phải bằng chứng an toàn. Lỗi tràn
bộ đệm trong cùng một struct allocation có thể tồn tại mà không sinh ra thông
báo lỗi nào. Đây là lý do tại sao các nguyên tắc lập trình an toàn
(\textit{snprintf}, kiểm tra độ dài, ownership rõ ràng) phải được áp dụng
\emph{ngay từ khi viết code}, không phải chỉ khi công cụ báo lỗi.
